\documentclass{scrreprt}
\usepackage{fhmsegu}
\usepackage{longtable}
\usepackage{booktabs}
\usepackage[all, defaultlines=5]{nowidow}
\usepackage{import}
\usepackage{setspace}

\addbibresource{quellen_methodik.bib}

\DeclareSIUnit{\euro}{\text{\euro}}
\DeclareSIUnit{\year}{a}
\DeclareSIUnit{\kwh}{kWh}
\DeclareSIUnit{\mwh}{MWh}

\makenomenclature
\graphicspath{{images/}}

\setlength{\parindent}{1.5em}
\setlength{\parskip}{1.0em}

\pgfplotsset{compat=1.18}
\pagestyle{scrheadings}
\clearpairofpagestyles

\ihead{\headmark}
\ohead{\getfhmsauthor}
\cfoot{\pagemark}
\pagestyle{scrheadings}
\setheadsepline{0.4pt}
\renewcommand*{\chapterpagestyle}{scrheadings}

\begin{document}

%Titeldefinition
\fhmstitle{Energetische Bewertung virtueller Gebäude anhand von Jahresdauerlinien der Lüftungswärmeverluste nach DIN V 18599}
\fhmssubtitle{Hausarbeit im Fach Energiepotentiale}
\fhmsauthor{Jan Kruse}
\fhmsdate{\today}
\fhmsemail{jan.kruse@kampmann.de}

\fhmsmaketitlepage
\pagestyle{empty}
\clearpage

\chapter*{Abstract}
% TODO: Kurzfassung DE + EN nach Fertigstellung

\pagestyle{scrheadings}
\pagenumbering{gobble}
\begin{spacing}{1.1}
\tableofcontents
\end{spacing}
\clearpage

\pagenumbering{Roman}
\setcounter{page}{0}
\listoffigures
\addcontentsline{toc}{chapter}{Abbildungsverzeichnis}
\listoftables
\addcontentsline{toc}{chapter}{Tabellenverzeichnis}

\par\vspace{1.0em}
\phantomsection
\addcontentsline{toc}{chapter}{Nomenklatur}
\markboth{Nomenklatur}{Nomenklatur}
{\parindent0pt
 \usekomafont{chapter}\usekomafont{chapterprefix}%
 \textsf{\textbf{Nomenklatur}}\par
}
\vspace{0.6em}
{
\renewcommand{\arraystretch}{1.4}
\begin{tabular}{ll}
    $\varPhi_{\mathrm{V}}$ & Lüftungswärmeverlust in \si{\kilo\watt} \\
    $\dot q_{\mathrm{v}}$ & Außenluftvolumenstrom in \si{\cubic\meter\per\hour} \\
    $\theta_{\mathrm{i}}$ & Raum-Solltemperatur in \si{\celsius} \\
    $\theta_{\mathrm{e}}$ & Außenlufttemperatur in \si{\celsius} \\
    $\theta_{\mathrm{v,mech}}$ & Zulufttemperatur-Sollwert der RLT-Anlage in \si{\celsius} \\
    $\rho$ & Dichte der Luft in \si{\kilo\gram\per\cubic\meter} \\
    $c_{\mathrm{p}}$ & spez. Wärmekapazität der Luft in \si{\kilo\joule\per\kilo\gram\per\kelvin} \\
    $Q_{\mathrm{V}}$ & Jahresenergiemenge Lüftung in \si{\kwh} \\
    $q_{\mathrm{h,18^\circ C,12h,mth}}$ & normierter monatlicher Heizenergiekennwert nach DIN V 18599-3 \\
    $g_{\mathrm{h,u/o,mth}}$ & Gradient des monatlichen Heizenergiekennwerts nach DIN V 18599-3 \\
    $f_{\mathrm{T,h}}$ & Korrekturfaktor Betriebszeit (Heizfall) nach DIN V 18599-3 \\
    $\theta_{\mathrm{hc}}$ & Zulufttemperatur-Sollwert im Kennwertverfahren nach DIN V 18599-3 in \si{\celsius} \\
    $t_{\mathrm{v,mech}}$ & tägliche Betriebszeit der RLT-Anlage nach DIN V 18599-3 in \si{\hour} \\
    $d_{\mathrm{v,mech}}$ & jährliche Betriebstage der RLT-Anlage nach DIN V 18599-3 in \si{\day\per\year} \\
    $\eta_{\mathrm{WRG}}$ & Temperaturänderungsgrad der Wärmerückgewinnung (-); entspricht $\eta_{\mathrm{t}}$ nach DIN V 18599-2 \\
    $\theta_{\mathrm{i,c}}$ & Raum-Solltemperatur Kühlung in \si{\celsius} \\
    $q_{\mathrm{c,18^\circ C,12h,mth}}$ & normierter monatlicher Kühlenergiekennwert nach DIN V 18599-3 \\
    $g_{\mathrm{c,u/o,mth}}$ & Gradient des monatlichen Kühlenergiekennwerts nach DIN V 18599-3 \\
    $f_{\mathrm{T,c}}$ & Korrekturfaktor Betriebszeit (Kühlfall) nach DIN V 18599-3 \\
    $x$ & Wasserdampfgehalt der Luft in \si{\gram\per\kilo\gram} \\
    $x_{\mathrm{soll}}$ & Mindest-Wasserdampfgehalt der Zuluft in \si{\gram\per\kilo\gram} \\
    $h$ & spezifische Enthalpie feuchter Luft in \si{\kilo\joule\per\kilo\gram} \\
    $P_{\mathrm{V}}$ & elektrische Ventilatorleistung in \si{\kilo\watt} \\
    $\Delta p_{\mathrm{ZUL/ABL}}$ & Druckerhöhung Zu-/Abluftventilator in \si{\pascal} \\
    $\eta_{\mathrm{V}}$ & Gesamtwirkungsgrad Ventilator/Motor/Antrieb (-) \\
    $f_{\mathrm{p}}$ & Druckverhältnis-Zahl (Teillast) nach DIN V 18599-3 (-) \\
    $W_{\mathrm{V}}$ & Jahresenergiemenge Ventilatoren in \si{\kwh} \\
    % TODO: weitere Symbole ergänzen
\end{tabular}
}

\addchap{Abkürzungsverzeichnis}
\begin{acronym}[ASHRAE]
    \acro{RLT}{Raumlufttechnik}
    \acro{NWG}{Nichtwohngebäude}
    \acro{TRY}{Testreferenzjahr}
    \acro{TRJ}{Testreferenzjahr-Datensatz}
    \acro{DWD}{Deutscher Wetterdienst}
    \acro{BBR}{Bundesamt für Bauwesen und Raumordnung}
    \acro{JDL}{Jahresdauerlinie}
    \acro{WRG}{Wärmerückgewinnung}
    \acro{SFP}{Specific Fan Power}
    \acro{VAV}{Variable Air Volume}
    % TODO: weitere Abkürzungen bei Einbau in Fließtext ergänzen
\end{acronym}

\pagestyle{scrheadings}

\chapter{Einleitung}
\pagenumbering{arabic}

Ziel dieser Arbeit ist die Entwicklung eines Python-basierten Berechnungstools, mit dem die energetischen Auswirkungen von Lüftungsanlagen für unterschiedliche Gebäudetypen und Standorte schnell abgeschätzt werden können. Hierzu werden auf Grundlage stündlicher Wetterdaten Jahresenergiemengen sowie Jahresdauerlinien für Heizung, Kühlung, Befeuchtung und Ventilatorbetrieb ermittelt. Die Ergebnisse werden anschließend exemplarisch mit dem Kennwertverfahren der DIN V 18599-3 verglichen
% Motivation: Anteil Lüftungswärmeverluste am Energiebedarf von NWG
% Ziel: Python-Tool zur überschlägigen Jahresenergiemenge Lüftung + JDL
% Aufbau der Arbeit (Verweis auf Kap. 2-5)

\chapter{Theoretischer Hintergrund}\label{chap:Theorie}

\section{Energetische Bewertung von Gebäuden nach DIN V 18599}
Die Vornormenreihe DIN V 18599 regelt die energetische Bewertung von
Gebäuden hinsichtlich des Nutz-, End- und Primärenergiebedarfs für
Heizung, Kühlung, Lüftung, Trinkwarmwasser und Beleuchtung
\parencite{DIN1859910}. Sie gliedert sich in zwölf aufeinander
aufbauende Teile: Teil 1 stellt das allgemeine Bilanzierungsverfahren
und die Verknüpfung der übrigen Teile bereit, Teil 2 behandelt den
Nutzenergiebedarf für Heizen und Kühlen von Gebäudezonen, Teil 3 den
Nutzenergiebedarf für die energetische Luftaufbereitung, Teil 7 den
Endenergiebedarf von \ac{RLT}- und Klimakältesystemen für den
Nichtwohnungsbau sowie Teil 10 die Nutzungsrandbedingungen und
Klimadaten \parencite{DIN1859910}. Die energetische Bilanzierung
erfolgt dabei grundsätzlich monatsweise auf Basis eines
Referenzklimas für Deutschland \parencite{DIN185993}.

\section{Das Kennwertverfahren nach DIN V 18599-3}
Für die Ermittlung des Nutzenergiebedarfs raumlufttechnischer Anlagen
stellt DIN V 18599-3 ein sogenanntes Kennwertverfahren bereit
\parencite{DIN185993}. Dieses folgt dem Grundgedanken für die in der Praxis häufigsten \ac{RLT}-Systeme eine umfangreiche Variantenmatrix aus 46 Anlagenvarianten aufzustellen, die
sich hinsichtlich Feuchteanforderung, Befeuchtertyp,
\ac{WRG}-Typ und \ac{WRG}-Größe unterscheiden \parencite{DIN185993}.
Für jede dieser Varianten wurde der Nutzenergiebedarf vorab in einer
stündlichen Simulation für einen festgelegten Basisfall (Referenzklima,
Zulufttemperatur-Sollwert $\theta_{\mathrm{v,mech}}=\SI{18}{\celsius}$,
tägliche Betriebszeit $t_{\mathrm{v,mech}}=\SI{12}{\hour}$) ermittelt und
als normierter, monatlicher Energiekennwert (z.\,B.
$q_{\mathrm{h,18^\circ C,12h,mth}}$) tabellarisch hinterlegt
\parencite{DIN185993}.

Zur Anwendung auf ein konkretes Gebäude wird aus der Variantenmatrix
die zutreffende Anlagenvariante ausgewählt und der zugehörige
Kennwert mithilfe von Gradienten (für abweichende
Zulufttemperatur-Sollwerte) sowie Korrekturfaktoren (für abweichende
tägliche Betriebszeiten) auf die tatsächlichen Randbedingungen
umgerechnet und auf den tatsächlichen Volumenstrom denormiert
\parencite{DIN185993}. Der Wärmerückgewinnungsgrad einer Anlage kann
dabei durch lineare Interpolation zwischen den in der Variantenmatrix
hinterlegten Stützstellen (45\,\%, 60\,\%, 75\,\%) berücksichtigt
werden \parencite{DIN185993}. Das Verfahren erlaubt damit eine
monatsweise, überschlägige Bilanzierung, ohne dass für jedes Gebäude
eine vollständige stündliche Simulation der Anlagentechnik
durchgeführt werden muss. Alternativ lässt die Norm für nicht durch
die Variantenmatrix abgebildete Anlagenkonzepte auch ein direktes
Stundenschrittverfahren zu \parencite{DIN185993}.

\section{Testreferenzjahre als Wetterdatengrundlage}
Für die energetische Bewertung von Gebäuden und ihrer technischen
Anlagen werden repräsentative, stündlich aufgelöste Wetterdatensätze
benötigt. Diese liefern für ein vollständiges Jahr klimatologische Randbedingungen, ohne die
Extremwerte eines einzelnen realen Messjahres widerzuspiegeln
\parencite{DWDBBR2017}. Während frühere Testreferenzjahre auf 15
Klimaregionen mit jeweils einer Repräsentanzstation beruhten
\parencite[vgl.][]{DIN1859910}, wurden diese im Rahmen der
TRY-Weiterentwicklung durch ortsgenaue Datensätze mit einer
räumlichen Auflösung von etwa \SI{1}{\square\kilo\meter} für ganz
Deutschland abgelöst \parencite{Spekat2011,DWDBBR2017}. Die
Interpolation der Basisgrößen (u.\,a. Lufttemperatur,
Windgeschwindigkeit, kurz- und langwellige Strahlung) auf dieses
Zielraster erfolgt unter Verwendung von Stations-, Satelliten- und
Modelldaten, u.\,a. des Regionalklimamodells COSMO-CLM
\parencite{Krähenmann2016,DWDBBR2017}.

Ein sogenanntes \ac{TRY} wird dabei nicht künstlich erzeugt, sondern
aus realen, mehrtägigen Witterungsabschnitten des Zeitraums 1995 bis
2012 zusammengesetzt, die so ausgewählt werden, dass Mittelwert und
Streuung der stündlichen Lufttemperatur sowie der Mittelwert der
Globalstrahlung im Vergleich zum langjährigen Jahresgang
bestmöglich getroffen werden \parencite{DWDBBR2017}. Neben den \acp{TRY} werden zusätzlich extreme \acp{TRY} erstellt, welche Jahre mit extrem warmer beziehungsweise kalter Witterungsbedingungen abbilden. Außerdem werden prognostizierte \acp{TRY} auf Basis regionaler
Klimaprojektionen für den Zeitraum 2031 bis 2060 bereitgestellt
\parencite{DWDBBR2017}. Für die vorliegende Arbeit wird ausschließlich
das gegenwärtige \ac{TRY} für den Standort des Technologie Campus Steinfurt verwendet, da eine überschlägige
Jahresenergiemenge unter heutigen, repräsentativen
Witterungsbedingungen ermittelt werden soll. Der Temperatur-und Feuchteverlauf des Datensatzes werden in Kapite \ref{chap:Ergebnisse} zur Einordnung dargestellt.

\chapter{Methodik}\label{chap:Methodik}

Im Gegensatz zur normativen Berechnung nach DIN V 18599 steht nicht die Erstellung eines normkonformen Energiebedarfsnachweises im Vordergrund. Ziel ist vielmehr die Entwicklung eines einfach anwendbaren Werkzeugs zur schnellen energetischen Abschätzung von Lüftungsanlagen. Daher werden verschiedene Vereinfachungen getroffen, die eine kurze Rechenzeit und eine einfache Parametrierung ermöglichen. Ausgehend von stündlichen
Wetterdaten und den Nutzungsrandbedingungen des Gebäudes werden
nacheinander die Betriebszustände, die stündlichen Energieströme für
Heizen, Kühlen, Befeuchten und Ventilatoren, deren Jahressummen sowie
die \ac{JDL} bestimmt. In dieser Arbeit wird die nur Leistung einer Dampfbefeuchtung dargestellt. Zwar wird in der Regel bei adiabatischen Befeuchtern keine thermische Energie zugeführt, jedoch muss durch die zu erwartende Temperaturabsenkung eine Nacherwärmung vorgesehen werden. Außerdem könnte es durch die Verdungstungskühlung die Jahreskühlenergie beeinflussen. Die Berechnung des Temperaturabfalls ist in CoolProp zwar möglich wurde im Rahmen dieser Arbeit allerdings nicht durchgeführt. Die folgenden Abschnitte gehen den Ablauf des Programmes schrittweise durch.

\section{Schritt 1: Eingangsdaten}
Grundlage der Berechnung sind zwei Arten von Eingangsdaten: stündlich
aufgelöste Wetterdaten und die Nutzungsrandbedingungen des betrachteten
Gebäudes.

Jeder Datensatz enthält 8760 Stundenwerte meteorologischer Größen für ein Jahr \parencite{DWDBBR2017}, von denen für die Berechnung die Außenlufttemperatur $\theta_{\mathrm{e}}$, der Wasserdampfgehalt $x_{\mathrm{e}}$ und der Luftdruck der Außenluft ausgewertet werden.

Die Nutzungsrandbedingungen des Gebäudes orientieren sich an den
Nutzungsprofilen für \ac{NWG} nach DIN V 18599-10
\parencite{DIN1859910}. Je Gebäude werden festgelegt: der
Außenluftvolumenstrom $\dot q_{\mathrm{v}}$, die Raum-Solltemperatur
Heizen $\theta_{\mathrm{i}}$, wahlweise eine Raum-Solltemperatur Kühlen
$\theta_{\mathrm{i,c}}$, die täglichen Betriebszeiten, der
Wärmerückgewinnungsgrad $\eta_{\mathrm{WRG}}$ sowie
ein Mindest-Wasserdampfgehalt der Zuluft $x_{\mathrm{soll}}$ und die
für die Ventilatorberechnung benötigten Druckerhöhungen und der
Ventilatorwirkungsgrad $\eta_{\mathrm{V}}$. Die konkreten Zahlenwerte
werden den Tabellen der DIN V 18599-10 entnommen bzw. für die
untersuchten Varianten projektspezifisch festgelegt.

\section{Schritt 2: Bestimmung der Betriebszustände}
Für jede der 8760 Stunden des Jahres wird zunächst geprüft, ob sich das
Gebäude innerhalb seiner hinterlegten Betriebszeit befindet.
Außerhalb der Betriebszeit wird die \ac{RLT}-Anlage als außer Betrieb angenommen. Sämtliche im Folgenden beschriebenen Energieströme
sind für diese Stunden null, unabhängig von der Außentemperatur oder
-feuchte. In diesem Ansatz werden vereinfacht konstante Volumenströme in der \ac{RLT}-Anlage 

\section{Schritt 3: Stündliche Heizleistung}
Die energetische Bewertung von Gebäuden nach DIN V 18599 erfolgt im
Regelfall über ein monatliches Bilanzverfahren \parencite{DIN1859910};
für die Lüftungswärme- und -kältebedarfe raumlufttechnischer Anlagen
stellt DIN V 18599-3 hierzu das in Abschnitt~\ref{chap:Theorie}
beschriebene Kennwertverfahren bereit, das auf einer Variantenmatrix
von 46 RLT-Anlagen-Varianten mit monatlich normierten Energiekennwerten
sowie Korrekturfaktoren für Betriebszeit, Zulufttemperatur und
Wärmerückgewinnungsgrad beruht \parencite{DIN185993}.

Für die vorliegende, überschlägige Betrachtung wird stattdessen eine
vereinfachte stündliche Energiebilanz auf Basis der \ac{TRY}-Daten
angesetzt. Diese folgt dem Grundprinzip der Lüftungswärmeverlustbilanz
nach DIN V 18599-2 \parencite{DIN185992}, wird jedoch nicht mit
Monatsmittelwerten, sondern mit stündlichen Temperaturwerten
ausgewertet, um jährliche Energiebedarf abschätzen zu können. Der
Lüftungswärmeverlust einer Betriebsstunde $t$ mit
$\theta_{\mathrm{e}}(t) < \theta_{\mathrm{i}}$ berechnet sich zu

\begin{equation}\label{eq:QdotV}
    \dot Q_{\mathrm{V}}(t) = \rho \, c_{\mathrm{p}} \, \dot q_{\mathrm{v}}(t) \, \Delta\theta_{\mathrm{eff}}(t)
\end{equation}

mit der wirksamen Temperaturdifferenz

\begin{equation}\label{eq:dtheta}
    \Delta\theta_{\mathrm{eff}}(t) = (1 - \eta_{\mathrm{WRG}}) \, \bigl(\theta_{\mathrm{i}} - \theta_{\mathrm{e}}(t)\bigr).
\end{equation}

Außerhalb der Betriebszeit sowie für Stunden mit
$\theta_{\mathrm{e}}(t) \geq \theta_{\mathrm{i}}$ liegt kein Heizfall
vor, sodass $\dot Q_{\mathrm{V}}(t) = 0$ gesetzt wird. Die
Jahresenergiemenge Lüftung $Q_{\mathrm{V}}$ ergibt sich als Summe der
8760 Stundenwerte nach Gleichung \eqref{eq:QdotV}.

Gleichung \eqref{eq:dtheta} ist die unmittelbare Umformung von
DIN V 18599-2:2018-09, Abschnitt~6.3.3.5, Gleichung~(98), welche die
Zulufttemperatur $\theta_{\mathrm{v,mech}}$ bei einfachen
Lüftungssystemen mit ungeregeltem Wärmeübertrager und einer
Ablufttemperatur gleich der Innentemperatur definiert als
$\theta_{\mathrm{v,mech}} = \theta_{\mathrm{e}} + \eta_{\mathrm{t}} \,
(\theta_{\mathrm{i}} - \theta_{\mathrm{e}})$
\parencite{DIN185992}. Löst man diese Gleichung nach der Differenz
zwischen Raum-Solltemperatur und Zulufttemperatur auf, ergibt sich
$\theta_{\mathrm{i}} - \theta_{\mathrm{v,mech}} = (1 -
\eta_{\mathrm{t}}) \, (\theta_{\mathrm{i}} - \theta_{\mathrm{e}})$,
also genau die in Gleichung \eqref{eq:dtheta} verwendete wirksame
Temperaturdifferenz $\Delta\theta_{\mathrm{eff}}$ mit
$\eta_{\mathrm{t}} = \eta_{\mathrm{WRG}}$.

Diese Vereinfachung unterscheidet sich in zwei wesentlichen Punkten
vom Kennwertverfahren nach DIN V 18599-3: Erstens erfolgt die
Bilanzierung stündlich statt monatlich, wodurch auf eine gesonderte
Umrechnung über Gradienten $g_{\mathrm{h,u/o,mth}}$
\parencite{DIN185993} verzichtet werden kann, da die tatsächliche
Zulufttemperatur stündlich in die Bilanz eingeht. Zweitens wird der
Einfluss der täglichen Betriebszeit nicht über einen empirischen
Korrekturfaktor $f_{\mathrm{T,h}}$ \parencite{DIN185993} abgebildet,
sondern unmittelbar über die stundengenaue Bestimmung der
Betriebszustände (Schritt~2) berücksichtigt.

\section{Schritt 4: Stündliche Kühlleistung}
Sofern für das virtuelle Gebäude eine Kühl-Solltemperatur
$\theta_{\mathrm{i,c}}$ hinterlegt ist, wird analog zum Heizfall ein
sensibler Kühlfall berechnet. Für Betriebsstunden mit
$\theta_{\mathrm{e}}(t) > \theta_{\mathrm{i,c}}$ gilt:

\begin{equation}\label{eq:QdotC}
    \dot Q_{\mathrm{C}}(t) = \rho \, c_{\mathrm{p}} \, \dot q_{\mathrm{v}}(t) \, (1-\eta_{\mathrm{WRG}}) \, \bigl(\theta_{\mathrm{e}}(t) - \theta_{\mathrm{i,c}}\bigr).
\end{equation}

Außerhalb der Betriebszeit sowie für Stunden mit
$\theta_{\mathrm{e}}(t) \leq \theta_{\mathrm{i,c}}$ ist $\dot
Q_{\mathrm{C}}(t) = 0$. Ist für ein Gebäude keine Kühl-Solltemperatur
festgelegt, entfällt dieser Schritt vollständig und die Kühlleistung
wird für alle Stunden zu null angenommen.

\section{Schritt 5: Stündliche Befeuchtungsleistung}
Sofern für das virtuelle Gebäude ein Mindest-Wasserdampfgehalt der
Zuluft $x_{\mathrm{soll}}$ hinterlegt ist, wird zusätzlich eine latente
Befeuchtungsbilanz gebildet. Für Betriebsstunden mit
$x_{\mathrm{e}}(t) < x_{\mathrm{soll}}$ gilt:

\begin{equation}\label{eq:Qst}
    \dot Q_{\mathrm{st}}(t) = \dot m_{\mathrm{tL}} \, \bigl(h(\theta_{\mathrm{i}}, x_{\mathrm{soll}}) - h(\theta_{\mathrm{i}}, x_{\mathrm{e}}(t))\bigr),
\end{equation}

wobei $h$ die spezifische Enthalpie feuchter Luft je kg trockener Luft
bezeichnet. Außerhalb der Betriebszeit sowie für Stunden mit
$x_{\mathrm{e}}(t) \geq x_{\mathrm{soll}}$ (die Außenluft ist bereits
feucht genug) ist $\dot Q_{\mathrm{st}}(t) = 0$. Der Wasserdampfgehalt
der Außenluft $x_{\mathrm{e}}(t)$ wird direkt der \ac{TRY}-Datei
entnommen; die Enthalpiewerte werden mit der Stoffdatenbibliothek
CoolProp \parencite{Bell2014} berechnet, um belastbare Stoffwerte
feuchter Luft ohne eigene, fehleranfällige Psychrometrie-Näherungsformeln
zu erhalten. Ist für ein Gebäude kein Feuchte-Sollwert festgelegt,
entfällt dieser Schritt und die Befeuchtungsleistung wird für alle
Stunden zu null angenommen.

\section{Schritt 6: Stündliche Ventilatorleistung}
Neben dem Nutzenergiebedarf für Heizen, Kühlen und Befeuchten wird
zusätzlich der elektrische Endenergiebedarf der Zu- und
Abluftventilatoren abgeschätzt, sofern für das Gebäude Druckerhöhungen
und ein Ventilatorwirkungsgrad hinterlegt sind. Nach DIN V 18599-3
berechnet sich die elektrische Ventilatorleistung allgemein aus dem
geförderten Volumenstrom, der Druckerhöhung, dem Ventilatorwirkungsgrad
sowie einer Druckverhältniszahl $f_{\mathrm{p}}$, die den Einfluss
einer Volumenstromregelung (\ac{VAV}) bei Teillast abbildet
\parencite{DIN185993}. Für Anlagen mit konstantem Volumenstrom, wie sie
den in dieser Arbeit betrachteten virtuellen Gebäuden zugrunde liegen,
entfällt die Teillastregelung ($f_{\mathrm{p}}=1$), sodass sich die
Gleichung auf den bekannten Sonderfall

\begin{equation}\label{eq:PV}
    P_{\mathrm{V}}(t) = \dot q_{\mathrm{v}}(t) \, \frac{\Delta p_{\mathrm{ZUL}} + \Delta p_{\mathrm{ABL}}}{\eta_{\mathrm{V}}}
\end{equation}

für Betriebsstunden reduziert \parencite{DIN185993}. Außerhalb der
Betriebszeit ist $P_{\mathrm{V}}(t) = 0$. Der Wirkungsgrad
$\eta_{\mathrm{V}}$ fasst dabei, analog zum \ac{SFP}-Wert nach
DIN V 18599-7, Abschnitt~3.1.2 (spezifische elektrische
Leistungsaufnahme je Volumenstrom), Ventilator-, Motor- und ggf.
Antriebswirkungsgrad zusammen \parencite{DIN185997}. Der Wirkungsgrad bezieht sich hier immer auf Zu- und Abluftventilatoren. Die Druckerhöhungen und Wirkungsgrade

Die in den Schritten 3 bis 6 beschriebenen Erweiterungen sind für eine
überschlägige Betrachtung vertretbar, vernachlässigen jedoch weiterhin
anlagenspezifische Regelungseffekte sowie im Fall der Befeuchtung die Unterscheidung zwischen Dampf- und Verdunstungsbefeuchtung, die im
vollständigen Kennwertverfahren über die Variantenmatrix erfasst werden
\parencite{DIN185993}.

\section{Schritt 7: Jahresenergiemenge und Volllaststunden}
Für jede der vier Energiearten (Heizen, Kühlen, Befeuchten,
Ventilatoren) wird die Jahresenergiemenge als Summe der 8760
Stundenwerte gebildet, z.\,B. für die Heizenergie

\begin{equation}
    Q_{\mathrm{V}} = \sum_{t=1}^{8760} \dot Q_{\mathrm{V}}(t).
\end{equation}

Zusätzlich werden die rechnerischen Volllaststunden als Quotient aus
Jahresenergiemenge und Spitzenlast bestimmt; sie geben an, wie viele
Stunden die Anlage bei Spitzenlast betrieben werden müsste, um
dieselbe Jahresenergiemenge zu erreichen.

\section{Schritt 8: Erstellung der Jahresdauerlinie}
Zur Erstellung der \ac{JDL} werden die 8760 Stundenwerte einer
Energieart absteigend sortiert. Die Abszisse der \ac{JDL} gibt dabei
die Anzahl der Jahresstunden an, in denen der jeweilige Leistungswert
erreicht oder überschritten wird. Dies erlaubt eine Einschätzung, wie
oft hohe Lasten im Jahresverlauf auftreten, und dient u.\,a. als
Grundlage für die Dimensionierung der Anlagentechnik. In der Folgenden Darstellung wird der Programmablauf als Flussdiagramm dargestellt.

\begin{figure}[H]
\centering
\subimport{TeX/}{Flussdiagramm_methodik.tex}
\caption{Flussdiagramm mit Programmablauf}
\label{abb:Flussdiagramm}
\end{figure}


\chapter{Ergebnisse \& Diskussion}\label{chap:Ergebnisse}

\section{Ergebnisse von Stichproben der virtuellen Gebäude}

Zur Bewertung des entwickelten Berechnungstools werden verschiedene exemplarische Gebäudeprofile untersucht. Die Auswahl umfasst typische Nichtwohngebäude mit unterschiedlichen Anforderungen hinsichtlich Betriebszeiten, Außenluftvolumenströmen, Kühlung, Wärmerückgewinnung und Befeuchtung. Die zugrunde gelegten Randbedingungen sind in Tabelle \ref{tab:ProfileServerraumLagerhalleGrossraumbuero} dargestellt. Die Parameter für die Gebäudeprofile werden als JSON Dateien im Projektordner hinterlegt. Es können beliebige Gebäudeprofile hinzugefügt werden.

\begin{table}[H]
    \centering
    \caption{Gebäudeprofile (JSON-Parameter des Tools)}
    \label{tab:ProfileServerraumLagerhalleGrossraumbuero}
    \begin{tabular}{lccc}
        \toprule
        Parameter & Serverräume & Lagerhalle & Bürogebäude \\
        \midrule
        $q_V$ [\si{\cubic\metre\per\hour}]& 2500 & 8000 & 5000 \\
        $\theta_{\text{i,h}}$ & 21,0 & 12,0 & 21,0 \\
        Betriebsstunden [h] & 24, Mo-So & 14, Mo-So & 11, Mo-Fr \\
        $\eta_{\text{WRG}}$ [-]& 0,6 & 0,0 & 0,6 \\
        $\theta_{\text{i,c}}$ [\si{\celsius}]& 24,0 & 26,0 & 24,0 \\
        $x_{\text{soll}}$ [\si{\gram\per\kilo\gram}]& - & - & 6,0 \\
        $\Delta P_{\text{ZUL}}$ [\si{\pascal}] & 500 & 400 & 350 \\
        $\Delta P_{\text{ABL}}$ [\si{\pascal}]& 400 & 300 & 280 \\
        $\eta_{\text{V}}$ & 0,65 & 0,55 & 0,60 \\
        \bottomrule
    \end{tabular}
\end{table}

\begin{figure}
\centering
\includegraphics[width=\textwidth]{TRY2015_Temperatur_Feuchte.pdf}
\caption{Temperatur- und Feuchteverlauf Standort Steinfurt aus \ac{TRY}}
\label{abb:Temperaturrelfeuchte}
\end{figure}

\section{Sensitivitätsbetrachtung}
% z. B. Einfluss Luftwechsel, Standort, WRG-Grad

\section{Vergleichsrechnung mit dem Kennwertverfahren nach DIN V 18599-3}\label{sec:Vergleichsrechnung}
Um die Ergebnisse des stündlichen Bilanzverfahrens einordnen zu
können, wird für zwei der virtuellen Gebäude zusätzlich eine
Vergleichsrechnung nach dem normativen Kennwertverfahren der DIN V
18599-3 (Jahresverfahren für Anlagen mit konstantem Volumenstrom,
Abschnitt~4.4.1 i.\,V.\,m. Tabelle~A.1) durchgeführt
\parencite{DIN185993}. Für das Einzelbüro-Profil wird dabei eine
Anlagenvariante ohne Feuchteanforderung herangezogen, die die
Funktionen Heizen und Kühlen abdeckt; für das Serverraum-Profil wird
zusätzlich eine Anlagenvariante mit Feuchteanforderung herangezogen,
sodass mit den beiden Vergleichsrechnungen alle vier vom Tool
berechneten Energiearten (Heizen, Kühlen, Befeuchten, Ventilatoren)
gegen das normative Verfahren abgeglichen werden. Diese
Gegenüberstellung erfolgt exemplarisch für zwei Gebäude und zwei
Anlagenvarianten und erhebt keinen Anspruch auf eine vollständige
Anwendung der gesamten Variantenmatrix.

\subsection{Einzelbüro-Profil: Heizen und Kühlen}
Aus der Variantenmatrix nach DIN V 18599-3, Tabelle~7, wird zunächst
die zum Einzelbüro-Profil passende Anlagenvariante identifiziert. Das
Profil sieht keine Feuchteanforderung sowie einen
Wärmerückgewinnungsgrad von $\eta_{\mathrm{WRG}}=0{,}6$ vor; dies
entspricht der Variante~3 (keine Feuchteanforderung,
\ac{WRG}-Typ „nur Wärme“, Wärmerückgewinnungsgrad $60\,\%$)
\parencite{DIN185993}. Für diese Variante weist Tabelle~A.1
(Gesamtjahr) einen Heizenergiekennwert von
$q_{\mathrm{h,18°C,12h}}=\SI{1148}{\watt\hour\hour\per\cubic\meter}$
sowie die Gradienten $g_{\mathrm{h,u}}=\SI{274}{\watt\hour\hour\per\kelvin\per\cubic\meter}$
(für $\theta_{\mathrm{hc}}<\SI{18}{\celsius}$) und
$g_{\mathrm{h,o}}=\SI{783}{\watt\hour\hour\per\kelvin\per\cubic\meter}$
(für $\theta_{\mathrm{hc}}>\SI{18}{\celsius}$) aus
\parencite{DIN185993}.

Für das Einzelbüro-Profil ($\theta_{\mathrm{hc}}=\SI{20}{\celsius}$,
tägliche Betriebszeit $t_{\mathrm{v,mech}}=\SI{11}{\hour}$
[7--18\,Uhr], $d_{\mathrm{v,mech}}=\SI{260}{\day\per\year}$
[Mo--Fr], $\dot q_{\mathrm{v}}=\SI{500}{\cubic\meter\per\hour}$)
ergibt sich der Heizenergiekennwert nach Umrechnung auf die
abweichende Zulufttemperatur zu

\begin{equation}
    q_{\mathrm{h}}(\theta_{\mathrm{hc}}) = q_{\mathrm{h,18°C,12h}} + g_{\mathrm{h,o}} \, (\theta_{\mathrm{hc}} - \SI{18}{\celsius}) \approx \SI{2714}{\watt\hour\hour\per\cubic\meter}
\end{equation}

Die Umrechnung auf die tatsächliche tägliche Betriebszeit erfolgt
über den Korrekturfaktor $f_{\mathrm{T,h}}$, der für
$t_{\mathrm{v,mech}}=\SI{11}{\hour}$ nahe $1{,}0$ liegt
\parencite{DIN185993}. Nach Skalierung auf die tägliche Betriebszeit,
die tatsächliche Anzahl jährlicher Betriebstage (Tabelle~A.1 basiert
auf $365$~Tagen) sowie Denormierung auf den Außenluftvolumenstrom von
$\SI{500}{\cubic\meter\per\hour}$ ergibt sich ein jährlicher
Nutzenergiebedarf Heizen von rund $Q_{\mathrm{h,Kennwert}}\approx
\SI{880}{\kwh\per\year}$. Diese Berechnung wurde zusätzlich
programmtechnisch nachgebildet und gegen den von Hand nachvollzogenen
Wert abgesichert.

Analog wird für den Kühlfall der Kühlenergiekennwert
$q_{\mathrm{c,18°C,12h}}=\SI{2309}{\watt\hour\hour\per\cubic\meter}$
sowie die Gradienten $g_{\mathrm{c,u}}=\SI{856}{\watt\hour\hour\per\kelvin\per\cubic\meter}$
und $g_{\mathrm{c,o}}=\SI{389}{\watt\hour\hour\per\kelvin\per\cubic\meter}$
der Variante~3 herangezogen \parencite{DIN185993}. Die reale
Raum-Kühlsolltemperatur des Einzelbüro-Profils nach DIN V 18599-10
beträgt $\theta_{\mathrm{i,c}}=\SI{24}{\celsius}$
\parencite{DIN1859910}; dieser Wert liegt jedoch außerhalb des in
Gleichung~(31)/(32) der DIN V 18599-3 festgelegten
Gültigkeitsbereichs von $\SI{14}{\celsius}$ bis $\SI{22}{\celsius}$
für den Zulufttemperatur-Sollwert \parencite{DIN185993}. Um eine
unzulässige Extrapolation zu vermeiden, wird für die
Vergleichsrechnung ersatzweise der obere Rand des Gültigkeitsbereichs
($\SI{22}{\celsius}$) angesetzt. Mit dem zugehörigen Korrekturfaktor
$f_{\mathrm{T,c}}$ für $t_{\mathrm{v,mech}}=\SI{11}{\hour}$
\parencite{DIN185993} ergibt sich nach gleichem Rechengang wie beim
Heizfall ein jährlicher Nutzenergiebedarf Kühlen von rund
$Q_{\mathrm{c,Kennwert}}\approx\SI{249}{\kwh\per\year}$. Für die
Befeuchtung ist innerhalb der Variante~3 kein Vergleich möglich, da
diese Variante keine Feuchteanforderung vorsieht und der zugehörige
Dampfkennwert $q_{\mathrm{st,12h}}=0$ ist \parencite{DIN185993}.

\subsection{Serverraum-Profil: Heizen, Kühlen und Befeuchten}
Damit auch die im Tool implementierte Befeuchtungsfunktion (Schritt~5)
gegen das Kennwertverfahren abgeglichen werden kann, wird zusätzlich
eine Vergleichsrechnung für das Serverraum-Profil durchgeführt. Dieses
Profil sieht neben Heizen und Kühlen einen Mindest-Wasserdampfgehalt
der Zuluft von $x_{\mathrm{soll}}=\SI{6.0}{\gram\per\kilo\gram}$ sowie
einen Wärmerückgewinnungsgrad von $\eta_{\mathrm{WRG}}=0{,}6$ vor. Nach
Tabelle~C.1 der DIN V 18599-3 entspricht ein Zuluft-Wasserdampfgehalt
von $\SI{6.0}{\gram\per\kilo\gram}$ genau dem Standardwert für
Feuchteanforderungen „mit Toleranzbereich“ \parencite{DIN185993}; das
Profil ist damit dieser Klassifizierung zuzuordnen. In Kombination mit
einem Dampfbefeuchter sowie dem \ac{WRG}-Typ „nur Wärme“
($60\,\%$) ergibt sich nach Tabelle~7 die Variante~21
\parencite{DIN185993}. Für diese Variante weist Tabelle~A.1
(Gesamtjahr) folgende Kennwerte aus:
$q_{\mathrm{h,18°C,12h}}=\SI{1028}{\watt\hour\hour\per\cubic\meter}$,
$g_{\mathrm{h,u}}=\SI{227}{\watt\hour\hour\per\kelvin\per\cubic\meter}$,
$g_{\mathrm{h,o}}=\SI{882}{\watt\hour\hour\per\kelvin\per\cubic\meter}$,
$q_{\mathrm{st,18°C,12h}}=\SI{3992}{\watt\hour\hour\per\cubic\meter}$,
$q_{\mathrm{c,18°C,12h}}=\SI{2443}{\watt\hour\hour\per\cubic\meter}$,
$g_{\mathrm{c,u}}=\SI{886}{\watt\hour\hour\per\kelvin\per\cubic\meter}$
sowie
$g_{\mathrm{c,o}}=\SI{246}{\watt\hour\hour\per\kelvin\per\cubic\meter}$
\parencite{DIN185993}.

Das Serverraum-Profil ist mit $\theta_{\mathrm{hc}}=\SI{21}{\celsius}$,
$\dot q_{\mathrm{v}}=\SI{1500}{\cubic\meter\per\hour}$ sowie
durchgehendem Betrieb ($t_{\mathrm{v,mech}}=\SI{24}{\hour}$,
$d_{\mathrm{v,mech}}=\SI{365}{\day\per\year}$) hinterlegt. Die
Umrechnung des Heizenergiekennwerts auf die abweichende
Zulufttemperatur ergibt
$q_{\mathrm{h}}(\theta_{\mathrm{hc}}) = q_{\mathrm{h,18°C,12h}} +
g_{\mathrm{h,o}} \, (\theta_{\mathrm{hc}} - \SI{18}{\celsius}) \approx
\SI{3674}{\watt\hour\hour\per\cubic\meter}$. Da die tägliche
Betriebszeit mit $\SI{24}{\hour}$ deutlich von der Basiszeit von
$\SI{12}{\hour}$ abweicht, fällt der Korrekturfaktor $f_{\mathrm{T,h}}$
nach Gleichung~(37) der DIN V 18599-3 mit $f_{\mathrm{T,h}}\approx
1{,}06$ merklich von $1{,}0$ ab \parencite{DIN185993}. Nach Skalierung
auf die tatsächliche Betriebszeit sowie Denormierung auf den
Außenluftvolumenstrom ergibt sich ein jährlicher Nutzenergiebedarf
Heizen von rund $Q_{\mathrm{h,Kennwert}}\approx
\SI{11700}{\kwh\per\year}$.

Für den Kühlfall wird, analog zum Einzelbüro-Profil, die reale
Raum-Kühlsolltemperatur des Serverraum-Profils
($\theta_{\mathrm{i,c}}=\SI{24}{\celsius}$) durch den oberen Rand des
Gültigkeitsbereichs der Umrechnungsgleichungen
($\theta_{\mathrm{hc}}=\SI{22}{\celsius}$) ersetzt, um eine unzulässige
Extrapolation zu vermeiden \parencite{DIN185993}. Es ergibt sich
$q_{\mathrm{c}}(\theta_{\mathrm{hc}}) = q_{\mathrm{c,18°C,12h}} -
g_{\mathrm{c,o}} \, (\theta_{\mathrm{hc}} - \SI{18}{\celsius}) \approx
\SI{1459}{\watt\hour\hour\per\cubic\meter}$. Mit dem für
Feuchteanforderungen „mit Toleranzbereich“ geltenden Korrekturfaktor
$f_{\mathrm{T,c}}$ nach Gleichung~(39) der DIN V 18599-3, der für
$t_{\mathrm{v,mech}}=\SI{24}{\hour}$ mit $f_{\mathrm{T,c}}\approx
0{,}66$ deutlich unter $1{,}0$ liegt \parencite{DIN185993}, ergibt sich
nach gleichem Rechengang ein jährlicher Nutzenergiebedarf Kühlen von
rund $Q_{\mathrm{c,Kennwert}}\approx\SI{2890}{\kwh\per\year}$.

Für die Dampfbefeuchtung gilt näherungsweise
$q_{\mathrm{st}}(\theta_{\mathrm{hc}}) \approx q_{\mathrm{st,18°C,12h}}$
unabhängig von der Zulufttemperatur \parencite{DIN185993}; der
zugehörige Korrekturfaktor $f_{\mathrm{T,st}}$ ist konstant $1{,}0$
\parencite{DIN185993}. Nach Skalierung auf die tatsächliche
Betriebszeit sowie Denormierung auf den Außenluftvolumenstrom ergibt
sich ein jährlicher Nutzenergiebedarf Befeuchtung von rund
$Q_{\mathrm{st,Kennwert}}\approx\SI{12000}{\kwh\per\year}$. Mit dieser
zweiten Vergleichsrechnung sind alle vier vom Tool berechneten
Energiearten – Heizen, Kühlen, Befeuchten und Ventilatoren, Letztere
über den bereits in Schritt~6 unmittelbar aus Gleichung~(11)/(12) der
DIN V 18599-3 abgeleiteten Sonderfall konstanten Volumenstroms – gegen
das normative Kennwertverfahren nachvollzogen.

\section{Einordnung und Grenzen der überschlägigen Methode}
Der Vergleich des stündlichen Bilanzverfahrens mit dem
Kennwertverfahren nach DIN V 18599-3 (Abschnitt~\ref{sec:Vergleichsrechnung})
zeigt, in welcher Größenordnung sich die überschlägige Methode
gegenüber dem normativen Verfahren bewegt. % TODO: Vergleich mit tatsächlichem TRY-Tool-Ergebnis ergänzen, sobald ein reales TRY für den betrachteten Standort ausgewertet wurde.

Bei der Interpretation der Abweichung ist zu berücksichtigen, dass
die in DIN V 18599-3, Tabelle~A.1 hinterlegten Energiekennwerte
selbst auf Grundlage eines Testreferenzjahres für die
„Region~4“ des früheren, auf 15 Klimaregionen basierenden
Referenzklimasystems nach DIN V 18599-10, Anhang~E ermittelt wurden
\parencite{DIN185993,DIN1859910}. Das im Rahmen dieser Arbeit
verwendete Berechnungstool greift hingegen auf die seit 2017
verfügbaren, ortsgenauen \acp{TRY} zurück
\parencite{DWDBBR2017,Spekat2011}. Eine Abweichung zwischen beiden
Verfahren ist daher nicht ausschließlich auf die methodische
Vereinfachung der stündlichen Bilanz zurückzuführen, sondern auch auf
eine unterschiedliche Wetterdatengrundlage. Für eine methodisch
sauberere Gegenüberstellung müssten beide Verfahren auf denselben
Klimadatensatz angewendet werden.

Darüber hinaus vernachlässigt das im Tool implementierte
Bilanzverfahren, wie in Abschnitt~\ref{chap:Methodik} dargelegt, den
latenten Wärmeanteil sowie anlagenspezifische Regelungseffekte, die
im Kennwertverfahren über die Variantenmatrix und die zugehörigen
Korrekturfaktoren implizit erfasst werden. Die überschlägige Methode
ist damit als Orientierungsrechnung zu verstehen, die für einen
schnellen Vergleich mehrerer Gebäudevarianten und Standorte sowie zur
Erstellung einer \ac{JDL} geeignet ist, jedoch keinen Ersatz für eine
normkonforme Nachweisführung nach DIN V 18599 darstellt.

Für den Kühlfall kommt als zusätzliche Einschränkung hinzu, dass die
reale Raum-Kühlsolltemperatur beider betrachteter Gebäudeprofile
($\SI{24}{\celsius}$) außerhalb des Gültigkeitsbereichs der
Umrechnungsgleichungen des Kennwertverfahrens liegt
(Abschnitt~\ref{sec:Vergleichsrechnung}); die Vergleichsrechnungen
Kühlen sind daher nur eingeschränkt aussagekräftig. Für die
Befeuchtung konnte anhand des Serverraum-Profils zwar ein
normkonformer Abgleich mit einer Anlagenvariante mit
Feuchteanforderung durchgeführt werden; die im Tool implementierte
latente Bilanz (Schritt~5) verwendet dabei jedoch reale, stündlich
schwankende Wasserdampfgehalte der Außenluft anstelle eines fest
vorgegebenen Standardwertes, sodass auch hier eine methodische
Vereinfachung gegenüber dem Kennwertverfahren besteht.

Für die berechnete Jahresenergiemenge der Ventilatoren gilt
einschränkend, dass die verwendeten Druckerhöhungen und der
Ventilatorwirkungsgrad projektspezifisch angenommen und nicht aus
DIN V 18599-7 entnommen wurden, da die Norm hierfür primär auf
anlagenspezifische Bestandsdaten (Inspektionswerte) verweist
\parencite{DIN185997}. Zudem wird mit $f_{\mathrm{p}}=1$ nur der
Sonderfall konstanten Volumenstroms abgebildet; für Anlagen mit
\ac{VAV}-Teillastregelung wäre die vollständige Gleichung nach
DIN V 18599-3 heranzuziehen \parencite{DIN185993}.

\chapter{Fazit \& Ausblick}\label{chap:Fazit}

\printbibliography

\appendix
\chapter{Parameterübersicht virtuelle Gebäude}
\chapter{Programmübersicht Python-Tool}

\end{document}
